% !TeX root = ../../Book3.tex
\chapter{整函数的增长}
\label{b3:ch:19}
\BookChapterNumber{b3:ch:19}
\begin{chapterabstract}
整函数在无穷远的大小可以从圆周最大模、Taylor 系数和零点计数三个方面观察。本章先用三圆定理建立对数半径上的凸性，并比较最大模与最大单项；随后定义增长阶与型，证明系数判据及导数、坐标变化和多项式复合下的增长规律。Jensen 公式把增长界转成零点收敛指数与密度的限制。最后在扇形和条带中用增长条件控制无穷远边界，分别建立严格次临界和临界最小型的 Phragmén--Lindelöf 比较，以及允许竖向增长的三线估计。
\end{chapterabstract}

\begin{learninggoals}
\begin{itemize}
\item 区分圆周最大模与平均对数，用三圆不等式及其等号情形识别径向增长的限制。
\item 比较最大模与 Taylor 最大项，追踪主导次数随半径的变化，并由全局幂增长判别多项式。
\item 正确使用阶与型中的上极限，通过系数衰减计算或构造给定增长，保留零阶、最小型和最大型的边界情形。
\item 说明导数、仿射变换与多项式复合如何改变阶型，识别加法与乘法中可能出现的抵消。
\item 按重数计算零点收敛指数，区分指数、临界幂和与密度常数提供的不同资料。
\item 在无界区域比较中逐一控制有限边界、顶点与远处边界，辨明严格指数条件、临界最小型条件和条带宽度的作用。
\end{itemize}
\end{learninggoals}

多项式的次数可以由 \(M(r)\) 的幂增长读出，而 \(e^{z^2}\) 的最大模是 \(e^{r^2}\)，必须先取对数才能看见幂次。另一方面，\(e^{\mathrm iz^2}\) 在实轴和虚轴上都具有模一值，却在两轴之间的某些方向快速增长。一个径向尺度如何从所有方向中取得，几条射线上的信息又在什么条件下足以约束整个平面，是本章始终需要分清的两个问题。

\ProseParagraphBreak
\secref{b3:sec:1901}先建立最大模的基本结构。全局幂界结合 Cauchy 系数估计使高次系数消失；三圆定理则用环域上的调和函数插值两条圆周的模界。等号情形还要利用复函数的刚性，才能从径向直线恢复单项式。Taylor 最大项提供另一个可算的径向量，其对数是一族整数斜率直线的上包络，但比较它与整个级数时仍须控制其余各项。

\ProseParagraphBreak
\secref{b3:sec:1902}把这些估计组织成阶与型。阶确定 \(\log M(r)\) 的幂尺度，型比较这个尺度下的系数；零阶仍容许超越整函数，给定正阶也容许零型和最大型。系数判据的两个方向分别来自最优选择 Cauchy 半径和控制模级数尾部。读者可以先跟算基本例子，再用这些判据处理稀疏系数；指标中有大量零系数，并不会使非零子列的上极限失去作用。

\ProseParagraphBreak
\secref{b3:sec:1903}回到零点。收敛指数可以由幂和或整数计数的增长读出，临界指数处的幂和则还取决于更细的修正。Jensen 公式使最大模限制零点数量，并把给定密度转成型的必要下界。无零点因子能够改变增长，所以这些比较的方向和前提应当保留；仅知道一个计数上界，不能反推出函数的最大模上界。

\ProseParagraphBreak
\secref{b3:sec:1904}处理无界区域的边界比较。严格次临界增长可以由更高次的衰减指数压住；达到临界指数时，这个余量消失，需要最小型条件及另一种调和屏障。有限多条射线把平面分割后，每个扇形都须满足自己的角宽条件。条带中的三线比较同样需要控制横向截断边界，允许的竖向增长与条带宽度相联系。边界模的比较与连续零边界的局部唯一性承担不同任务，不应把增长条件不加区分地附到所有唯一性问题上。

\ProseParagraphBreak
本章使用 Cauchy 公式与系数估计、次调和比较，以及\chapref{b3:ch:18}的 Jensen 公式和零点计数。Levin 的《Lectures on Entire Functions》\cite{Levin1996Lectures}可作为继续学习增长方向与零点分布的参考。本章先建立后续整函数乘积理论所需的增长尺度；一般指标理论和给定零点的乘积构造不作为本章证明的前提。

% !TeX root = ../../Book3.tex
\section{最大模}
\label{b3:sec:1901}

多项式与指数函数都在全平面全纯，却有不同的无穷远尺度。最大模把方向各异的增长压成一个径向量，三圆定理则限制这个量怎样随半径变化。

\subsection{\texorpdfstring{\(M(r,f)\)}{M(r,f)}}
\label{b3:subsec:190101}

对整函数 \(f\)，定义最大模 \(M(r,f)=\max_{|z|=r}|f(z)|\)。最大模原理使它等于闭圆盘上的最大模，并随 \(r\) 非减。若 \(f\) 非常值，则 \(M\) 严格递增：两半径相同最大值会在大圆盘内部取得最大模，迫使函数常数。

多项式 \(a_dz^d+\cdots+a_0\) 有 \(M(r)=|a_d|r^d(1+o(1))\)；\(e^z\) 有 \(M(r)=e^r\)。沿负实轴的指数趋零，与它的最大模指数增长并不矛盾。

\subsection{最大模}
\label{b3:subsec:190102}

\begin{proposition}\label{b3:prop:19-polynomial-growth}
若整函数满足 \(M(r,f)\le C(1+r)^N\)，则为次数不超过 \(\lfloor N\rfloor\) 的多项式。
\end{proposition}
\begin{proof}
Cauchy 估计给 \(|f^{(k)}(0)|\le k!C(1+r)^N/r^k\)。当整数 \(k>N\) 时令 \(r\to\infty\)，得到导数零；Taylor 展开因而有限。反向显然由有限项三角不等式成立。
\end{proof}
这里控制的是所有方向，而不是一条射线。函数 \(e^{-z}\) 在正实轴有界，却不是多项式。

\subsection{Hadamard 三圆定理}
\label{b3:subsec:190103}

\term[hadamard-sanyuan]{Hadamard 三圆定理}[Hadamard three-circles theorem]可由次调和比较直接得到。
\begin{theorem}[Hadamard 三圆]\label{b3:thm:19-three-circles}
若 \(f\) 在环域 \(r_1\le|z|\le r_3\) 邻域全纯且不恒零，则 \(r_1<r_2<r_3\) 时
\[
\log M(r_2)\le
\frac{\log(r_3/r_2)}{\log(r_3/r_1)}\log M(r_1)
+\frac{\log(r_2/r_1)}{\log(r_3/r_1)}\log M(r_3).
\]
\end{theorem}
\begin{proof}
选择 \(A,B\) 使调和函数 \(A\log|z|+B\) 在两边界等于相应最大对数。函数 \(\log|f|-A\log|z|-B\) 次调和，两边界均非正，最大值原理给内部非正。在半径 \(r_2\) 取最大值，线性插值就是题式。
\end{proof}

\subsection{对数凸性}
\label{b3:subsec:190104}

令 \(t=\log r\)，三圆定理说 \(t\mapsto\log M(e^t)\) 凸。其割线斜率随半径向外不减；这是增长约束，而不是通常变量 \(r\) 下的凸性。

例如单项式 \(z^d\) 给严格直线 \(dt\)；指数给 \(e^t\)。如果某非零整函数在所有大半径的对数斜率都不超过整数 \(d\)，则 \(M(r)\le Cr^d\)，上一命题将它识别为多项式。这使一个径向几何性质决定整个函数的代数结构。


\ProseParagraphBreak

三圆定理也能在同一半径上比较几何平均与最大值。Jensen 的圆周平均描述零点总影响，最大模则取最有利增长的方向；前者可能恒定，后者却增长很快。对 \(e^z\)，圆周对数平均零，而 \(\log M(r)=r\)。因此由平均对数的凸性不能直接推出最大模的精确增长，三圆定理另需对边界最大值作调和比较。这两种凸性虽然都以 \(\log r\) 为变量，控制的对象仍不同。

\ProseParagraphBreak

最大模关于半径连续。若 \(r_j\to r>0\)，在共同紧环带上 \(f\) 一致连续，将 \(r_je^{\mathrm{i}t}\) 与 \(re^{\mathrm{i}t}\) 比较，函数值差一致趋零，故两圆周最大值之差趋零。对非恒零整函数 \(M(r)>0\)，对数也连续。三圆不等式因此给一个通常意义的连续凸函数，而非只在排除零点的半径上定义的对象。圆周即使碰到函数零点，最大模仍由其他位置取得，不需要移开半径。

\ProseParagraphBreak

若三圆不等式在一个中间半径取得等号，比较函数 \(\log|f|-A\log|z|-B\) 在环域内部取得最大值零，故恒为零。于是 \(|f(z)|=e^B|z|^A\)，\(f\) 在环域无零点。局部对数微分给 \(f'/f=A/z\)，沿绕原点的圆周积分，辐角原理使 \(A\) 必为整数；继而 \(f(z)=cz^A\) 于环域。若 \(f\) 还在原点全纯，整数 \(A\) 必非负。三圆等号的刚性来自次调和强最大值原理，不能只由一个实凸函数的一次等号就辨认出复函数的形式。

\ProseParagraphBreak

最大模还可与 Taylor 展开的最大单项比较。对 \(f(z)=\sum c_nz^n\)，定义 \(\mu(r)=\max_n|c_n|r^n\)。整性使项随 \(n\to\infty\) 趋零，所以最大值确实达到。Cauchy 系数界给 \(\mu(r)\le M(r)\)；另一方面，对 \(0<r<R\)，
\[
 M(r)\le\sum_n|c_n|r^n
       \le\mu(R)\sum_n(r/R)^n
       =\frac{\mu(R)}{1-r/R}.
\]
这两边估计说明，模级数中的相消可能影响某个方向的函数值，却不能使整个圆周上的最大模比所有单项都小。取 \(R=2r\)，两种径向量具有相同的增长阶；当讨论有限正阶的型时，再取 \(R=(1+\varepsilon)r\) 并让 \(\varepsilon\downarrow0\)，它们的型也相同。
\ProseParagraphBreak
从对数坐标看，\(\log\mu(e^t)=\max_n(\log|c_n|+nt)\)，是斜率为非负整数的一族直线的上包络，因而凸。取得最大值的指标随 \(t\) 非减：若在 \(t_1<t_2\) 分别由 \(n_1,n_2\) 取到最大值，把两次比较不等式相加便得 \((n_2-n_1)(t_2-t_1)\ge0\)。这使 Taylor 展开的主导项随半径逐渐移向高次幂。若主导指标最终有共同上界，则所有足够高次系数都必须为零，否则该项的斜率更大，最终会超过已有的有限条直线。多项式与超越整函数的差别，由此再次表现为无穷远是否不断出现新的主导次数。

% !TeX root = ../../Book3.tex
\section{增长阶与型}
\label{b3:sec:1902}

多项式的次数可以从 \(\log M(r)/\log r\) 读出，但这个量对每个超越整函数都趋于无穷。若 \(M(r)\) 大致像 \(e^{r^s}\)，应当转而比较 \(\log M(r)\) 与 \(r^s\)。阶用来确定这里的幂指数，型则在同一幂尺度上保留前面的系数。先把这两层比较的量词说清楚，再研究它们怎样由导数、复合和 Taylor 系数决定。

\subsection{阶}
\label{b3:subsec:190201}

\Needspace{6\baselineskip}
\begin{definition}\label{b3:def:19-order}
设 \(f\) 为非常值整函数，\(M(r,f)=\max_{|z|=r}|f(z)|\)。其\term[zengzhang-jie]{增长阶}[order of growth]定义为
\[
 \rho(f)\coloneqq\limsup_{r\to\infty}
 \frac{\log\log M(r,f)}{\log r}\in[0,\infty].
\]
常数函数，包括零函数，约定增长阶为零。增长阶有限的整函数称为\term[youxianjie-zhenghanshu]{有限阶整函数}[entire function of finite order]。
\end{definition}
\symidx{\(\rho(f)\)：整函数的增长阶}
非常值时 \(M(r,f)\to\infty\)，所以定义中的双对数只需在充分大半径使用。其上极限不会为负：取某个 \(c_k\ne0\)、\(k\ge1\)，便有 \(M(r,f)\ge|c_k|r^k\)，从而双对数的商的下极限不小于零。采用上极限，是为了容许增长在不同半径上出现变化，不预先假定存在一个精确的渐近式。

\ProseParagraphBreak

阶最常用的含义是可用哪些指数上界。严格大于阶的指数总能给出上界；严格小于阶的指数则无论乘上多大的常数都不够。

\Needspace{6\baselineskip}
\begin{proposition}[阶与指数增长尺度]\label{b3:prop:19-growth-scales}
设 \(f\) 为整函数，其增长阶为 \(\rho\in[0,\infty]\)。对每个实数 \(s>\rho\)，存在 \(C_s>0\)，使
\[
 M(r,f)\le C_s e^{r^s},\quad r\ge0.
\]
对每个 \(0<s<\rho\) 及任意 \(A,C>0\)，存在任意大的半径 \(r\) 使 \(M(r,f)>C e^{Ar^s}\)。因此
\[
 \rho=\inf\bigl\{\,s>0\mid
 \exists A,C>0\ \forall r\ge0,\ M(r,f)\le C e^{Ar^s}\,\bigr\},
\]
其中空集的下确界约定为无穷。
\end{proposition}
\begin{proof}
常数函数直接满足结论，以下设 \(f\) 非常值。若 \(s>\rho\)，上极限的定义给充分大半径上的
\[
 \log\log M(r,f)\le s\log r,
 \quad\text{即}\quad \log M(r,f)\le r^s.
\]
再用 \(M\) 在有限半径区间的连续性扩大 \(C_s\)，得到全半径的估计。

若 \(0<s<\rho\)，选 \(s'\) 使 \(s<s'<\rho\)；\(\rho=\infty\) 时也可如此选择。存在趋于无穷的半径序列 \(r_j\)，使 \(\log\log M(r_j,f)>s'\log r_j\)。故
\[
 \frac{\log M(r_j,f)}{r_j^s}>r_j^{s'-s}\longrightarrow\infty.
\]
对给定的 \(A,C\)，充分大的 \(j\) 便满足 \(\log M(r_j,f)>A r_j^s+\log C\)。最后，任何上界 \(M(r,f)\le C e^{Ar^s}\) 都使双对数的上极限不超过 \(s\)；结合已经证明的严格阈值，得到下确界公式。
\end{proof}
下确界不一定在 \(s=\rho\) 处取得。知道 \(\rho<\infty\)，只保证每个 \(\varepsilon>0\) 都能给出 \(M(r)\le C_\varepsilon e^{r^{\rho+\varepsilon}}\)。能否在临界指数 \(\rho\) 上取得有限常数，需要再看型。

\subsection{型}
\label{b3:subsec:190202}

\Needspace{6\baselineskip}
\begin{definition}\label{b3:def:19-type}
设整函数 \(f\) 的增长阶为 \(\rho\in(0,\infty)\)，\(M(r,f)=\max_{|z|=r}|f(z)|\)。其在自身增长阶上的\term[zengzhang-xing]{型}[type]定义为
\[
 \tau(f)\coloneqq\limsup_{r\to\infty}
 \frac{\log M(r,f)}{r^\rho}\in[0,\infty].
\]
\(\tau=0\)、\(0<\tau<\infty\)、\(\tau=\infty\) 分别称为\term[zuixiaoxing]{最小型}[minimal type]、\term[zhengchangxing]{正常型}[normal type]、\term[zuidaxing]{最大型}[maximal type]。
\end{definition}
\symidx{\(\tau(f)\)：整函数在自身有限正阶上的型}
这里的尺度必须是函数自身的阶。在更大指数上，这个上极限自动为零；在更小的正指数上，它自动为无穷，均不能细分同阶函数。零阶不按上述公式作统一型分类：非常值多项式与超越零阶函数都有 \(\log M(r)\to\infty\)，除以 \(r^0=1\) 后无法区分它们。无穷阶也不属于此定义的范围。

\ProseParagraphBreak

与阶的阈值类似，若 \(\tau<\infty\)，则对每个 \(b>\tau\)，存在 \(C_b>0\) 使
\[
 M(r,f)\le C_b e^{b r^\rho},\quad r\ge0.
\]
若 \(0<b<\tau\)，则 \(M(r,f)e^{-b r^\rho}\) 沿某个趋于无穷的半径序列趋于无穷。前者来自上极限的最终上界，后者可选 \(b'\in(b,\tau)\)，在满足 \(\log M(r)>b'r^\rho\) 的序列上验证。因而 \(\tau\) 是所有能给出上述全半径上界的正数 \(b\) 的下确界；\(\tau=\infty\) 时这一集合为空。

\ProseParagraphBreak

临界常数本身仍未必可用。例如 \(f(z)=(1+z)e^z\) 满足
\[
 M(r,f)=(1+r)e^r.
\]
上界来自 \(|1+z|\le1+r\)、\(\operatorname{Re}z\le r\)，而在正实点 \(z=r\) 同时取得等号。于是 \(\rho=1\)、\(\tau=1\)，却不存在统一的 \(C\) 使 \(M(r,f)\le C e^r\)。因此从型的定义得到的 \(\tau+\varepsilon\) 通常不能直接换成 \(\tau\)。

\ProseParagraphBreak

前节的最大项也能读出这两个量。比较型时必须让两个半径的比可以任意接近一；只比较 \(r\) 和 \(2r\)，会留下一个无法消去的常数因子。

\Needspace{6\baselineskip}
\begin{proposition}[最大项给出相同的阶与型]\label{b3:prop:19-maximum-term-growth}
设 \(f(z)=\sum_{n=0}^{\infty}c_nz^n\) 为非常值整函数，\(\mu(r,f)=\max_{n\ge0}|c_n|r^n\)。则
\[
 \rho(f)=\limsup_{r\to\infty}
 \frac{\log\log\mu(r,f)}{\log r}.
\]
若 \(0<\rho(f)<\infty\)，则其型还满足
\[
 \tau(f)=\limsup_{r\to\infty}
 \frac{\log\mu(r,f)}{r^{\rho(f)}}.
\]
两个公式均容许其所定义的量为零或无穷。
\end{proposition}
\begin{proof}
某个正次系数非零，所以 \(\mu(r,f)\to\infty\)，两个公式在大半径有意义。固定 \(\kappa>1\)，由\cref{b3:prop:19-maximum-term}，
\[
 \mu(r,f)\le M(r,f)
 \le\frac{\kappa}{\kappa-1}\mu(\kappa r,f).
\]
乘以一个正常数只使双对数改变 \(o(1)\)，把半径 \(r\) 换成 \(\kappa r\) 则有 \(\log(\kappa r)/\log r\to1\)。两边取双对数的上极限，便得到阶相同，包括零阶和无穷阶。

设共同的阶为 \(\rho\in(0,\infty)\)，并把最大项公式右端记为 \(\tau_\mu\)。取一次对数并除以 \(r^\rho\)，上述不等式给
\[
 \tau_\mu\le\tau(f)\le\kappa^\rho\tau_\mu.
\]
若 \(\tau_\mu\) 有限，让 \(\kappa\downarrow1\) 得到相等；若 \(\tau_\mu=\infty\)，左侧已经迫使 \(\tau(f)=\infty\)。
\end{proof}

\subsection{基本例子}
\label{b3:subsec:190203}

先看 \(f(z)=e^{az^m}\)，其中 \(a\ne0\)、\(m\ge1\) 为整数。在 \(|z|=r\) 上有 \(\operatorname{Re}(az^m)\le|a|r^m\)，又能选取一个角度使 \(az^m\) 为正实数，因而
\[
 M(r,f)=e^{|a|r^m},\quad \rho(f)=m,\quad \tau(f)=|a|.
\]
这里必须选到使实部最大的位置；只看实轴，在 \(a\) 为纯虚数时便会遗漏全部增长。

\ProseParagraphBreak

三角函数也属于一阶正常型。由指数表达式，
\[
 |\sin z|\le\frac{e^{\operatorname{Im}z}+e^{-\operatorname{Im}z}}2\le e^r,
 \quad
 |\cos z|\le e^r,\quad |z|=r.
\]
在虚轴上则有 \(|\sin(\mathrm{i}r)|=\sinh r\)、\(|\cos(\mathrm{i}r)|=\cosh r\)。两者均为 \(\tfrac12e^r(1+o(1))\)，所以两种最大模都满足 \(\log M(r)=r+O(1)\)，得到阶一、型一。它们在实轴有界与此并不矛盾。

\ProseParagraphBreak

多项式和迭代指数给出两个端点。若 \(P\) 为 \(d\ge1\) 次多项式，\(\log M(r,P)=d\log r+O(1)\)，从而 \(\rho(P)=0\)，无论次数多大。另一方面，\(f(z)=e^{e^z}\) 满足
\[
 |f(z)|=e^{\operatorname{Re}(e^z)}\le e^{e^r},\quad |z|=r,
\]
在 \(z=r\) 取得等号。因此 \(\log\log M(r,f)=r\)，阶为无穷。

\ProseParagraphBreak

零阶还包含超越整函数。考虑
\[
 F(z)=\sum_{n=0}^{\infty}e^{-n^2}z^n.
\]
系数的 \(n\) 次根为 \(e^{-n}\)，趋于零，所以收敛半径无穷。所有系数为正，三角不等式的上界在 \(z=r\) 取得，于是令 \(t=\log r\) 后有
\[
 M(r,F)=F(r)
 =\sum_{n=0}^{\infty}e^{-n^2+nt}
 =e^{t^2/4}\sum_{n=0}^{\infty}e^{-(n-t/2)^2}.
\]
当 \(t\ge1\) 时，选离 \(t/2\) 最近的非负整数 \(n_0\)，其距离不超过 \(1/2\)，便得到最后一和的下界 \(e^{-1/4}\)。对其他指标按 \(|n-n_0|\) 分组，则该和至多为
\[
 1+2\sum_{k=1}^{\infty}e^{-(k-1/2)^2},
\]
这是与 \(t\) 无关的有限常数。因此
\[
 \log M(r,F)=\frac14(\log r)^2+O(1).
\]
它超过每个固定的 \(d\log r\)，却小于每个 \(r^\varepsilon\) 的最终尺度；所以 \(F\) 超越而阶为零。最大模的对数斜率仍趋于无穷，符合前节的凸斜率判别。

\ProseParagraphBreak

一般的代数运算先给上界。若 \(f,g\) 为整函数，则
\[
 M(r,f+g)\le M(r,f)+M(r,g),\quad
 M(r,fg)\le M(r,f)\cdot M(r,g).
\]
结合指数尺度判据可得
\[
 \rho(f+g)\le\max\{\rho(f),\rho(g)\},\quad
 \rho(fg)\le\max\{\rho(f),\rho(g)\}.
\]
例如当两阶有限时，对每个大于二者的 \(s\)，和、积都满足 \(C e^{A r^s}\) 型上界，再让 \(s\) 下降即可；若最大阶无穷，上界本身自动成立。若两函数阶不同，和的阶必等于较大者，否则把高阶函数写成和减去低阶函数，将与刚才的上界矛盾。同阶相消则可能使阶下降：\(e^z+(-e^z+1)=1\)，而乘积 \(e^ze^{-z}=1\) 也说明乘法一般没有对应的等式。

\ProseParagraphBreak

平移和缩放改变圆心与半径，却不会任意改变增长尺度。用闭圆盘最大值表示 \(M\)，可以把这种变化写成精确的包含关系。

\Needspace{6\baselineskip}
\begin{proposition}[仿射变量变换]\label{b3:prop:19-affine-growth}
设 \(f\) 为整函数，\(a,b\in\mathbb C\)、\(a\ne0\)，并令 \(g(z)=f(az+b)\)。则
\[
 \rho(g)=\rho(f).
\]
若 \(\rho(f)=\rho\in(0,\infty)\)，则
\[
 \tau(g)=|a|^\rho\tau(f),
\]
其中型容许为零或无穷。
\end{proposition}
\begin{proof}
仿射映射把 \(|z|\le r\) 送到以 \(b\) 为中心、半径 \(|a|r\) 的闭圆盘。当 \(|a|r>|b|\) 时，这个圆盘包含以原点为中心、半径 \(|a|r-|b|\) 的闭圆盘，并包含于半径 \(|a|r+|b|\) 的闭圆盘。因此
\[
 M(|a|r-|b|,f)\le M(r,g)\le M(|a|r+|b|,f).
\]
常数函数的结论直接成立。若 \(f\) 非常值，令 \(s=|a|r\pm|b|\)，则每种代换都遍历一个充分大的半径区间，且 \(\log s/\log r\to1\)。对夹逼取双对数的上极限，得到阶不变，包括零阶和无穷阶。

若共同的阶为 \(\rho\in(0,\infty)\)，则
\[
 \frac{\log M(s,f)}{r^\rho}
 =\left(\frac{s}{r}\right)^\rho
   \frac{\log M(s,f)}{s^\rho},
 \quad \frac{s}{r}\longrightarrow|a|.
\]
夹逼的两端因此有同一上极限 \(|a|^\rho\tau(f)\)，得到型的公式。
\end{proof}
特别地，平移不改变阶与型。缩放的影响则保留在型的因子 \(|a|^\rho\) 中。

\ProseParagraphBreak

把仿射映射换成 \(d\) 次多项式，半径 \(r\) 的主尺度就变为 \(r^d\)。上界只需三角不等式；为了得到反向估计，还要保证多项式的像覆盖一个足够大的圆盘。

\Needspace{6\baselineskip}
\begin{proposition}[多项式复合的增长]\label{b3:prop:19-polynomial-composition}
设 \(f\) 为整函数，\(P(z)=a_dz^d+\cdots+a_0\) 为 \(d\ge1\) 次多项式，\(a_d\ne0\)。则
\[
 \rho(f\circ P)=d\rho(f),
\]
其中零阶和无穷阶按扩展数值解释。若 \(\rho(f)=\rho\in(0,\infty)\)，则在 \(f\circ P\) 自身的阶 \(d\rho\) 上，
\[
 \tau(f\circ P)=|a_d|^\rho\tau(f).
\]
\end{proposition}
\begin{proof}
固定 \(0<q<|a_d|<Q\)。当 \(r\) 充分大时，三角不等式给
\[
 |P(z)|\le Qr^d,\quad |z|\le r.
\]
另一方面，对任意 \(|w|\le q r^d\)，在 \(|z|=r\) 上都有
\[
 |P(z)-w-a_dz^d|
 \le\sum_{k<d}|a_k|r^k+q r^d
 <|a_d|r^d.
\]
\(P-w\) 与 \(a_dz^d\) 在闭圆盘邻域全纯，且满足严格的边界比较，故\cref{b3:thm:03-rouche}使二者在盘内有相同的零点数 \(d\)。特别地，每个这样的 \(w\) 都是某个 \(|z|<r\) 的像。于是
\[
 M(qr^d,f)\le M(r,f\circ P)\le M(Qr^d,f).
\]

若 \(f\) 为常数，结论直接成立。其余情形将 \(s=q r^d\) 或 \(s=Q r^d\) 代入，有 \(\log s/\log r\to d\)，而 \(s\) 遍历所有充分大的半径。夹逼两端的双对数商都以上极限 \(d\rho(f)\) 为界，故得到阶的等式。

设 \(0<\rho<\infty\)。同一夹逼除以 \(r^{d\rho}\) 后给
\[
 q^\rho\tau(f)\le\tau(f\circ P)\le Q^\rho\tau(f).
\]
令 \(q\uparrow|a_d|\)、\(Q\downarrow|a_d|\)，即得有限型的公式；若 \(\tau(f)=\infty\)，左边已经迫使复合函数也为无穷型，零型则由右边直接得到。
\end{proof}
因此，多项式的次数决定复合后阶的倍数，首项系数决定型的缩放，低次项只影响较小的半径误差。这一论证同时使用了像的上界与覆盖性质，单凭 \(|P(z)|\le C r^d\) 只能得到单向增长估计。

\subsection{导数增长与系数判据}
\label{b3:subsec:190204}

Cauchy 估计把较大圆周上的函数值转成较小圆周上的导数界；沿线段积分又能从导数恢复函数。二者合起来说明微分不会改变阶，在有限正阶上也不会改变型。

\Needspace{6\baselineskip}
\begin{proposition}[导数的阶与型]\label{b3:prop:19-derivative-order}
设 \(f\) 为整函数，常数函数的阶按零计。则
\[
 \rho(f')=\rho(f).
\]
若 \(\rho(f)=\rho\in(0,\infty)\)，则 \(\tau(f')=\tau(f)\)，包括最小型与最大型。
\end{proposition}
\begin{proof}
常数和一次多项式的结论直接成立，以下设 \(f'\) 非常值。对每个 \(|z|\le r\)，以 \(z\) 为圆心、\(r\) 为半径使用 Cauchy 导数估计，所用圆盘包含于 \(|\zeta|\le2r\)，故
\[
 M(r,f')\le\frac{M(2r,f)}r.
\]
取 \(r>1\)，双对数与固定半径倍增给 \(\rho(f')\le\rho(f)\)。反过来，沿从 \(0\) 到 \(z\) 的线段积分，得到
\[
 |f(z)|\le|f(0)|+|z|M(|z|,f'),
 \quad
 M(r,f)\le|f(0)|+rM(r,f').
\]
因 \(f'\) 非常值，某个正次 Taylor 系数非零，Cauchy 估计使 \(M(r,f')\ge c r^k\)，其中 \(c>0\)、\(k\ge1\)。于是 \(\log r=O(\log M(r,f'))\)，积分估计便给
\[
 \log\log M(r,f)\le\log\log M(r,f')+O(1).
\]
除以 \(\log r\) 并取上极限，得反向阶不等式。整个论证同样适用于零阶与无穷阶。

为比较型，固定 \(\varepsilon>0\)，在以 \(z\) 为中心、半径 \(\varepsilon r\) 的圆周上使用 Cauchy 估计，得到
\[
 M(r,f')\le\frac{M((1+\varepsilon)r,f)}{\varepsilon r}.
\]
取对数、除以 \(r^\rho\) 并取上极限，给出
\[
 \tau(f')\le(1+\varepsilon)^\rho\tau(f).
\]
积分估计在大半径可写为 \(M(r,f)\le2rM(r,f')\)，故 \(\tau(f)\le\tau(f')\)，因为 \(\log(2r)/r^\rho\to0\)。若 \(\tau(f)\) 有限，让 \(\varepsilon\downarrow0\) 得到相等；若 \(\tau(f)=\infty\)，反向不等式已经迫使 \(\tau(f')=\infty\)。
\end{proof}
反复使用命题可知，任何固定阶导数都保留原函数的阶；在有限正阶上还保留型。对整函数的任意全纯原函数应用同一命题，可知原函数的阶也相同，在有限正阶上型也相同。积分常数只改变零次项。

\ProseParagraphBreak

Taylor 系数则给出更直接的判据。Cauchy 估计允许为每个系数选择不同的半径：当最大模近似 \(e^{r^s}\) 时，最有利的半径由指数增长与分母 \(r^n\) 的竞争决定。反过来，从系数估计恢复最大模时，除了控制最大的单项，还必须控制有效项的数量和剩余尾部。

\Needspace{6\baselineskip}
\begin{theorem}[增长的系数判据]\label{b3:thm:19-coefficient-growth}
设 \(f(z)=\sum_{n=0}^{\infty}c_nz^n\) 为超越整函数。则
\[
 \rho(f)=\limsup_{\substack{n\to\infty\\c_n\ne0}}
 \frac{n\log n}{\log(1/|c_n|)}.
\]
非零系数的分母最终为正，公式容许零阶与无穷阶。若 \(\rho=\rho(f)\in(0,\infty)\)，其型还满足
\[
 \limsup_{n\to\infty}n^{1/\rho}|c_n|^{1/n}
 =(e\rho\tau(f))^{1/\rho},
\]
等价地，
\[
 e\rho\tau(f)=\limsup_{n\to\infty}n|c_n|^{\rho/n}.
\]
型的两个公式中，零系数对应的项取零，等式按零值和无穷值的意义解释。
\end{theorem}
\begin{proof}
先证明阶公式。整性使 \(|c_n|^{1/n}\to0\)，所以非零系数最终满足 \(|c_n|<1\)，分母为正。超越性又保证非零系数有无穷多个，因而指定子序列上的上极限有意义。记这个上极限为 \(a\in[0,\infty]\)。

若 \(\rho(f)<\infty\)，取任意 \(s>\rho(f)\)。指数尺度判据给 \(M(r,f)\le C e^{r^s}\)，Cauchy 系数界于是为
\[
 |c_n|\le C e^{r^s}r^{-n},\quad r>0.
\]
为了使右边尽量小，对 \(r^s-n\log r\) 求极小值。导数为零时 \(sr^s=n\)，代入 \(r=(n/s)^{1/s}\) 得
\[
 \log(1/|c_n|)
 \ge\frac ns\log n-\frac ns(1+\log s)-\log C.
\]
右边的主项为 \((n/s)\log n\)，其余为 \(O(n)\)，故 \(a\le s\)。让 \(s\downarrow\rho(f)\)，得到 \(a\le\rho(f)\)；若 \(\rho(f)=\infty\)，这一不等式自动成立。

反过来，若 \(a<\infty\)，取任意 \(s>a\)。上极限的定义使充分大的非零系数满足
\[
 \frac{n\log n}{\log(1/|c_n|)}<s,
 \quad\text{从而}\quad |c_n|\le n^{-n/s}.
\]
零系数也满足最后的界。将不满足此界的有限头部记为多项式 \(P\)，并只考虑 \(r\ge1\)。对尾部的单项，把指标看成正实数 \(x\)，其对数由
\[
 \phi(x)=x\log r-\frac{x}{s}\log x
\]
控制。\(\phi'(x)=0\) 给 \(x=r^s/e\)，且 \(\phi''(x)=-1/(sx)<0\)，故最大值为 \(r^s/(es)\)。在 \(n\le(2r)^s\) 的部分，项数至多 \(1+(2r)^s\)；在 \(n>(2r)^s\) 的部分，\(r n^{-1/s}<1/2\)，其和至多为 \(\sum_{n\ge1}2^{-n}=1\)。于是
\[
 M(r,f)\le M(r,P)+\bigl(1+(2r)^s\bigr)e^{r^s/(es)}+1.
\]
取对数后，有限多项式与项数至多贡献 \(O(\log r)\)，故 \(\log M(r,f)=O(r^s)\)，从而 \(\rho(f)\le s\)。让 \(s\downarrow a\) 得 \(\rho(f)\le a\)。当 \(a=0\) 时，此论证对每个 \(s>0\) 成立，给 \(\rho(f)=0\)；当 \(a=\infty\) 时，已证的 \(a\le\rho(f)\) 迫使 \(\rho(f)=\infty\)。阶公式的所有情形至此闭合。

下面固定 \(0<\rho=\rho(f)<\infty\)，并记
\[
 L=\limsup_{n\to\infty}n^{1/\rho}|c_n|^{1/n}.
\]
若 \(\tau(f)<\infty\)，任取 \(t>\tau(f)\)，有 \(M(r,f)\le C_t e^{t r^\rho}\)。最小化 \(t r^\rho-n\log r\)，应取 \(r^\rho=n/(\rho t)\)，于是
\[
 |c_n|\le C_t\left(\frac{e\rho t}{n}\right)^{n/\rho}.
\]
取 \(n\) 次根、乘 \(n^{1/\rho}\)，再取上极限，\(C_t^{1/n}\to1\)，所以 \(L\le(e\rho t)^{1/\rho}\)。令 \(t\downarrow\tau(f)\)，得到
\[
 L^\rho\le e\rho\tau(f).
\]
\(\tau(f)=0\) 时仍可从任意正 \(t\) 向下取极限；\(\tau(f)=\infty\) 时这个方向自动成立。

为证反向，若 \(L<\infty\)，任取 \(l>L\)。除有限项外，
\[
 |c_n|\le\bigl(l n^{-1/\rho}\bigr)^n.
\]
连续指标函数 \(x\log(lr)-(x/\rho)\log x\) 在 \(x=(lr)^\rho/e\) 取最大值 \((lr)^\rho/(e\rho)\)。将级数分在 \(n\le(2lr)^\rho\) 与 \(n>(2lr)^\rho\) 两部分：前一部分至多有 \(1+(2lr)^\rho\) 项，后一部分仍由几何级数 \(\sum2^{-n}\) 控制。连同有限头部，得到
\[
 \log M(r,f)\le\frac{l^\rho}{e\rho}r^\rho+O(\log r).
\]
除以 \(r^\rho\) 并取上极限，再令 \(l\downarrow L\)，便有 \(e\rho\tau(f)\le L^\rho\)。这也覆盖 \(L=0\)，因为可对每个 \(l>0\) 使用该估计。若 \(L=\infty\) 而 \(\tau(f)<\infty\)，与前一方向矛盾；若 \(\tau(f)=\infty\) 而 \(L<\infty\)，与后一方向矛盾。故两个端点同样相等。最后，正幂函数保持上极限，得到两种型公式的等价性。
\end{proof}
阶公式使用的是系数衰减的主尺度 \(n\log n\)，型则还保留其下一层指数常数。稀疏级数也可应用：零系数在阶公式中略去，在型公式中贡献零；但非零系数的下标必须始终使用它在原 Taylor 级数中的次数。

\ProseParagraphBreak

这些公式给出任意正实数阶的具体函数。设 \(s>0\)，令
\[
 F_s(z)=\sum_{n=0}^{\infty}\frac{z^n}{(n!)^{1/s}}.
\]
对递增函数 \(\log t\) 作积分比较，有
\[
 \int_1^n\log t\,\mathrm{d}t
 \le\log n!
 \le\int_1^n\log t\,\mathrm{d}t+\log n,
\]
因而 \(\log n!=n\log n-n+O(\log n)\)。特别地，\((n!)^{1/n}\sim n/e\)，系数根趋于零，故 \(F_s\) 为超越整函数。代入阶公式，
\[
 \frac{n\log n}{\log\bigl((n!)^{1/s}\bigr)}\longrightarrow s.
\]
再用型公式的第二种写法，
\[
 e s\tau(F_s)
 =\lim_{n\to\infty}\frac{n}{(n!)^{1/n}}=e,
\]
于是 \(\rho(F_s)=s\)、\(\tau(F_s)=1/s\)。例如 \(\sum z^n/(n!)^2\) 的阶为 \(1/2\)、型为 \(2\)。增长阶可以取非整数值，并不受多项式次数的整数性限制。

% !TeX root = ../../Book3.tex
\section{零点序列}
\label{b3:sec:1903}

Jensen 公式限制零点的数量，但还可以把整个零点序列压成一个收敛指数。下一章构造典范乘积时，这个指数决定必须扣掉多少低阶项。

\subsection{收敛指数}
\label{b3:subsec:190301}

按重数列出趋于无穷的非零零点 \(a_n\)。定义\term[lingdian-shoulian-zhishu]{零点收敛指数}[exponent of convergence]为 \(\lambda=\inf\{s>0:\sum|a_n|^{-s}<\infty\}\)。有限零点序列取零。对于无穷序列，若 \(n(r)\) 为计数，则
\[
\lambda=\limsup_{r\to\infty}\frac{\log n(r)}{\log r}.
\]
证明可用分部求和：\(\sum_{|a_n|>1}|a_n|^{-s}=s\int_1^\infty n(t)t^{-s-1}\,\mathrm{d}t\) 加有限边界项。若计数为 \(O(t^{s-\varepsilon})\) 积分收敛；若序列幂和收敛，则 \(n(r)\le r^s\sum|a_n|^{-s}\)，给反向。

\subsection{零点密度}
\label{b3:subsec:190302}

\begin{proposition}\label{b3:prop:19-zero-density}
有限阶 \(\rho\) 的非恒零整函数具有 \(n(r)=O_\varepsilon(r^{\rho+\varepsilon})\)，零点收敛指数不超过 \(\rho\)。
\end{proposition}
\begin{proof}
先除去原点零点。Jensen 估计取外半径 \(2r\)，得到 \(n(r)\log2\le\log M(2r)-\log|f(0)|\)。阶的定义给计数界；取稍小的 \(\varepsilon\) 后，上一积分判据使每个 \(s>\rho\) 的零点幂和收敛。
\end{proof}
原点的有限重数只增加常数项，不改变收敛指数。

\subsection{增长限制}
\label{b3:subsec:190303}

反向并不唯一：\(e^{z^m}\) 无零点而阶 \(m\)，所以零点只能约束增长的一部分。\(\sin\pi z\) 的零点为整数，收敛指数一，和阶一相同。

若指定 \(a_n=n^\alpha\)，其收敛指数为 \(1/\alpha\)。任何具有这些零点的整函数必须有阶至少 \(1/\alpha\)。下一章将证明离散零点总能实现，但若还要求有限阶，就要检查这个收敛指数。


\ProseParagraphBreak

零点计数的幂界可用二进尺度再证明收敛指数判据。若 \(n(r)\le Cr^a\)，在 \(2^j<|a_n|\le2^{j+1}\) 内，\(s>a\) 时幂和至多 \(C2^{(j+1)a-js}\)，对 \(j\) 求和为收敛几何级数。若 \(\sum|a_n|^{-s}<\infty\)，则半径 \(r\) 内每项至少 \(r^{-s}\)，故 \(n(r)r^{-s}\) 有界。这个分组证明不要求计数函数可微，也自然保留重数；分部求和则给更精确的常数与临界积分。

\ProseParagraphBreak

在临界指数上，下确界不决定幂和是否收敛。零点模长为 \(n\) 时收敛指数一，临界一次幂和发散；模长为 \(n\log^2(n+1)\) 时指数仍一，却临界一次幂和收敛。若模长为 \(n\log(n+1)\)，临界再次发散。对数修正不会改变 \(\log n(r)/\log r\) 的上极限，却决定下一章的最小初等因子次数。阶只决定粗尺度，属还要检查临界幂和，二者不可互换。

\ProseParagraphBreak

零点重数也能提高收敛指数。若每个正整数 \(n\) 的重数为 \(n\)，幂和变成 \(\sum n^{1-s}\)，收敛指数为二；计数约为 \(r^2/2\)，因此任何实现函数阶至少二。若仅按不同零点个数，则会得到错误下界一。对整函数增长、典范乘积和 Jensen 的贡献，统一按重数计数是必要条件，并非一个可选的记账方式。

\ProseParagraphBreak

增长下界只来自零点径向分布，角位置可以有进一步作用。例如整数零点在两条相反射线上排列，正弦给出阶一实现；把正整数零点全部放在同一射线，最小属仍是一，但沿某些方向的精确增长会不同。下一章的乘积修正会显式出现一次项，说明方向排列影响低次指数自由度。当前的 \(\lambda\le\rho\) 是普适必要条件，还没有声称它给所有方向的精确渐近。

\ProseParagraphBreak

若不仅知道阶，还知道正常型，Jensen 估计可给更具体的密度常数。假设 \(0<\rho<\infty\)、\(\tau<\infty\)，先去掉原点零点。对任意固定 \(u>1\)，有
\[
 n(r)\log u\le\log M(ur)-\log|f(0)|.
\]
除以 \(r^\rho\) 后取上极限，右侧不超过 \(\tau u^\rho\)。函数 \(u^\rho/\log u\) 在 \(\log u=1/\rho\) 处取最小值 \(e\rho\)，因此
\[
 \limsup_{r\to\infty}\frac{n(r)}{r^\rho}\le e\rho\tau.
\]
这仍是一般上界，不声称常数对每类函数都最优。它至少说明最小型 \(\tau=0\) 时 \(n(r)=o(r^\rho)\)，比只给 \(O(r^{\rho+\varepsilon})\) 更细。有限个原点零点除以 \(r^\rho\) 后消失，不改变结果。
\ProseParagraphBreak
对正弦函数，\(n(r)\) 与 \(r\) 同阶；对指数函数，计数恒为零，二者却有相同的阶与型。这再次说明增长和零点不是互相决定的两组资料。若只给定零点序列，下一章构造的典范乘积会固定其中一部分增长，而允许的无零点因子仍可改变最大模。使用密度估计时，正确方向是“增长上界限制零点”，不能不加额外假设把它倒置为“零点上界限制函数”。

% !TeX root = ../../Book3.tex
\section{Phragmén–Lindelöf 原理}
\label{b3:sec:1904}

最大值原理在无界扇形中不能直接应用。Phragmén--Lindelöf 方法先乘一个衰减指数，在大圆弧压住增长，再把有限区域的最大值原理送到无穷远。

\subsection{扇形域}
\label{b3:subsec:190401}

扇形 \(S=\{r e^{\mathrm{i}\theta}:r>0,\ |\theta|<\beta/2\}\)，\(0<\beta<2\pi\)，具有可选定的幂分支 \(z^p\)。若 \(p\beta/2<\pi/2\)，则 \(\operatorname{Re}z^p\ge c|z|^p\)，其中 \(c=\cos(p\beta/2)>0\)。这给整个闭扇形同方向的衰减，不能只检查中心射线。

\subsection{增长条件}
\label{b3:subsec:190402}

\term[phragmen-lindelof]{Phragmén--Lindelöf 原理}[Phragmén--Lindelöf principle]的一个常用版本如下。
\begin{theorem}[Phragmén--Lindelöf：扇形]\label{b3:thm:19-pl}
设 \(f\) 在上述扇形全纯并连续到有限边界与顶点。若两边界射线上 \(|f|\le1\)，且 \(|f(z)|\le C e^{A|z|^q}\)，其中 \(q<\pi/\beta\)，则 \(|f|\le1\) 于整个扇形。
\end{theorem}
\begin{proof}
选 \(q<p<\pi/\beta\)，令 \(f_\varepsilon=f e^{-\varepsilon z^p}\)。两射线上模不超过一；大圆弧上模不超过 \(C e^{Ar^q-\varepsilon c r^p}\)，充分大时亦不超过一。顶点值由连续性不超过一。对截断扇形用最大模原理得 \(|f_\varepsilon|\le1\)，再令大半径趋无穷，最后 \(\varepsilon\downarrow0\)。
\end{proof}
临界指数必须另设条件：\(e^{z^{\pi/\beta}}\) 在两边界模一，在中心射线却发散，给出严格不等式的必要性。

\subsection{整函数应用}
\label{b3:subsec:190403}

整函数若阶 \(\rho<1\)，并在实轴有界，则为常数。选择 \(\rho<q<1\)，阶的定义给 \(Ce^{Ar^q}\) 增长；分别对上、下半平面（角宽 \(\pi\)）应用扇形原理，得到全平面有界，再用 Liouville。

阶一的 \(e^{\mathrm{i}z}\) 在实轴模一而非常值，显示阈值不能只把严格不等号改成不大于。若函数仅在一个半平面全纯，结论只保证内部有界，不会自动成为常数。

\subsection{唯一性定理}
\label{b3:subsec:190404}

同一方法给唯一性。若 \(f,g\) 在扇形满足相同的次临界增长，且在两边界相等，则差 \(h\) 的边界零。对每个 \(\varepsilon>0\)，缩放 \(h/\varepsilon\) 满足边界模不超过一，扇形原理给 \(|h|\le\varepsilon\)，故 \(f=g\)。

这里边界相等不能直接当作恒等定理的内部聚点条件。本证明通过增长界得到全域比较；若边界为可作 Schwarz 反射的直线，还能用反射后的恒等定理得到更强的局部唯一性。临界例子中的零边值实部 \(\operatorname{Re}z^{\pi/\beta}\) 也说明无穷远会引入额外自由度。


\ProseParagraphBreak

扇形的旋转不改变临界阈值。若中轴辐角为 \(\theta_0\)，应在衰减因子中使用 \((e^{-\mathrm{i}\theta_0}z)^p\)，使实部在整个扇形有正下界；直接保留 \(z^p\) 可能在旋转后变成增长因子。幂的分支固定在该扇形中，连续延到两边界射线；顶点只有连续值，不必声称非整数幂在顶点全纯。截断扇形应用最大模原理时，顶点属于边界，内域不包含那个可能的分支点。

\ProseParagraphBreak

若 \(f\) 在顶点仅有界而未声明连续，证明可以先删去小圆。小圆上的模上界可能大于一，不能直接丢弃；若顶点为可去奇点，补上后再使用原定理即可。若允许顶点有极点，则边界射线在接近顶点仍受模一控制已经强烈限制它；但对于只控制远处边界的变式，小圆会留下额外常数。定理的有限边界与顶点假设用来控制截断域的全部边界，不能只关注大圆弧的增长。

\ProseParagraphBreak

一个多射线应用可以把阶阈值提高。若整函数阶 \(\rho\) 小于 \(\pi/\beta\)，并在若干从原点发出的射线上有共同模界 \(M\)，且相邻射线最大角宽为 \(\beta\)，则每个夹角扇形都可选增长指数 \(q\) 位于 \(\rho\) 与 \(\pi/\beta\) 之间。分别使用扇形定理，得到全平面 \(|f|\le M\)，故常数。例如四条坐标半轴把平面分成角宽 \(\pi/2\) 的扇形，于是阶小于二且在实轴、虚轴都有界 的整函数为常数。临界阶二的 \(e^{z^2}\) 在虚轴有界，在实轴却无界；\(e^{\mathrm{i}z^2}\) 在两轴均模一而非常值，给出正确的临界反例。

\ProseParagraphBreak

这里有限多个扇形的常数可以统一，因为只取有限次最大值；若射线无限多而各自界不同，不应未经估计就宣称有一个全平面共同界。把每条射线上的界与一个统一界区别开，也是使用 Liouville 结论前必须核查的量词。

\ProseParagraphBreak

边界零的唯一性应用需要同时假设两函数的增长。差函数的界由两者相加获得，常数可能加倍但指数仍是同一个次临界 \(q\)。然后对任意 \(\delta>0\) 使用 \(h/\delta\)，增长常数会随 \(\delta\) 改变；定理允许每次使用不同有限常数，因此最终 \(|h|\le\delta\) 合法。不能要求所有 \(\delta\) 共用一个增长常数，那会把证明中的无害缩放误认为失败。

\ProseParagraphBreak

无界区域比较还有一个常用的条带版本，它让边界模界不同的情形呈现出插值形式。
\begin{theorem}[有界函数的三线不等式]\label{b3:thm:19-three-lines}
设 \(f\) 在 \(0<\operatorname{Re}z<1\) 全纯，在闭条带连续有界，并满足 \(|f(\mathrm{i}t)|\le M_0\)、\(|f(1+\mathrm{i}t)|\le M_1\)，其中 \(M_0,M_1>0\)。则
\[
 |f(x+\mathrm{i}y)|\le M_0^{1-x}M_1^x\qquad(0\le x\le1).
\]
\end{theorem}
\begin{proof}
令 \(g(z)=f(z)\exp(-(1-z)\log M_0-z\log M_1)\)。实对数无分支问题，\(g\) 在闭条带有界，两竖边模不超过一。对 \(\varepsilon>0\)，令 \(g_\varepsilon(z)=g(z)e^{\varepsilon z^2-\varepsilon}\)。在两竖边有 \(\operatorname{Re}(\varepsilon z^2-\varepsilon)=\varepsilon(x^2-y^2-1)\le0\)。在横边 \(y=\pm T\)，其模不超过 \(\|g\|_\infty e^{-\varepsilon T^2}\)，充分大 \(T\) 时也不超过一。矩形最大模原理给 \(|g_\varepsilon|\le1\)，令 \(T\to\infty\)、\(\varepsilon\downarrow0\) 得 \(|g|\le1\)，还原即得结论。
\end{proof}
有界性在这条证明中控制横边，指数修正把横边送到无穷远后消去。若希望允许更快的竖向增长，必须选择新的修正并重新估计，不能直接引用当前版本。三线不等式的权重由观察点在两边界之间的相对位置决定，正如三圆定理的权重由对数半径决定；两个公式背后都是求出合适的调和比较函数。
\ProseParagraphBreak
如果某一边界模界为零，可先把该 \(M_j\) 替换为任意小正数，应用已经证明的正界版本再取极限；于是整个条带函数为零。这一处理不需要对零取对数，也不会留下未定义的指数因子。一般的边界参数极限应在不等式已经建立之后执行，而不是在构造辅助全纯函数时直接代入奇异参数。


\begin{chaptersummary}
三圆定理使 \(\log M(e^t)\) 成为凸函数，等号通过次调和强最大值原理限制全纯函数的形式。最大项的对数则是由 Taylor 系数决定的直线上包络。Cauchy 估计与模级数控制把这两个径向量联系起来，既给出多项式判别，也说明何时主导次数必须不断向高次移动。

\ProseParagraphBreak
阶与型记录不同层次的增长。阶为零仍可能是超越整函数，有限正阶也未必有有限正型。系数判据必须同时证明必要性和充分性，不能只由 Cauchy 估计得到单向限制。精确型的比较需要让半径比趋于一；固定倍增半径通常只足以确定阶。加法与乘法的上界允许抵消，具体等式还须检查函数之间的关系。

\ProseParagraphBreak
零点收敛指数由计入重数的幂和及计数增长共同刻画，下确界本身并不决定临界幂和是否收敛。Jensen 公式把零点的对数权重放在最大模之下，所以增长上界限制零点密度，既有密度又给增长提供必要下界。无零点的指数因子没有改变这些零点资料，却能改变型乃至阶。

\ProseParagraphBreak
无界区域的比较需要控制送往无穷远的那部分边界。扇形的临界指数由角宽决定，严格次临界情形保留衰减幂次的选择余量，临界最小型情形则用更精细的屏障。三线比较将两条边界的常数对数作调和插值，竖向增长条件保证截断边界的影响消失。把射线界或条带界用于全局结论时，应逐项检查这些增长条件及有限边界假设。
\end{chaptersummary}

\BookExercises

% 自拟：计算有限指数和的逐方向上极限；对并列主导指数用长区间平方平均排除持续相消，不调用一般指标理论。
\Needspace{6\baselineskip}
\begin{exercise}[有限指数和的方向增长]\label{b3:ex:19-draft-01}
设 \(\lambda_1,\ldots,\lambda_k\) 为互异复数，\(c_1,\ldots,c_k\ne0\)，且至少一个 \(\lambda_j\ne0\)。令
\[
 F(z)=\sum_{j=1}^kc_je^{\lambda_jz}.
\]
证明对每个固定 \(\theta\in\mathbb R\)，有
\[
 \limsup_{r\to\infty}\frac{\log|F(re^{\mathrm i\theta})|}{r}
 =\max_j\operatorname{Re}(\lambda_je^{\mathrm i\theta}).
\]
允许射线上出现零点，届时对数取 \(-\infty\)。特别要处理两个或更多指数具有相同最大实部的情形。由此求 \(F\) 的阶与型。
\end{exercise}
\begin{solution}
固定 \(\theta\)，记
\[
 H=\max_j\operatorname{Re}(\lambda_je^{\mathrm i\theta}),\quad
 J=\{j:\operatorname{Re}(\lambda_je^{\mathrm i\theta})=H\}.
\]
三角不等式给 \(|F(re^{\mathrm i\theta})|\le e^{Hr}\sum_j|c_j|\)，所以题中上极限不超过 \(H\)。为得反向估计，写
\[
 e^{-Hr}F(re^{\mathrm i\theta})=p(r)+E(r),\quad
 p(r)=\sum_{j\in J}c_je^{\mathrm i\omega_jr},\quad
 \omega_j=\operatorname{Im}(\lambda_je^{\mathrm i\theta}).
\]
不在 \(J\) 中的项具有严格较小实部，故 \(E(r)\to0\)；若没有这些项，取 \(E=0\)。\(J\) 中的实频率 \(\omega_j\) 互异，因为相同的实部和虚部会使两个 \(\lambda_j\) 相等。

将平方展开，在长区间 \([T,2T]\) 上平均，得到
\[
 \frac1T\int_T^{2T}|p(r)|^2\,\mathrm dr
 =\sum_{j\in J}|c_j|^2
 +\sum_{\substack{j,\ell\in J\\j\ne\ell}}
 c_j\overline{c_\ell}\frac1T\int_T^{2T}e^{\mathrm i(\omega_j-\omega_\ell)r}\,\mathrm dr.
\]
每个非对角积分除以 \(T\) 后的模不超过 \(2/(T|\omega_j-\omega_\ell|)\)，故平均趋于
\(S=\sum_{j\in J}|c_j|^2>0\)。因此对充分大的 \(T\)，区间内有一点 \(r_T\)，使 \(|p(r_T)|\ge\sqrt{S/2}\)。又因 \(E\) 在 \([T,\infty)\) 一致趋小，最终有
\[
 |F(r_Te^{\mathrm i\theta})|\ge\frac{\sqrt S}{4}e^{Hr_T}.
\]
沿 \(r_T\to\infty\) 取对数并除以 \(r_T\)，得到下界 \(H\)。这只需要找到不断远离原点的一列非相消位置，并未声称射线上存在普通极限。

令 \(L=\max_j|\lambda_j|>0\)。全局上界为 \(M(r,F)\le(\sum|c_j|)e^{Lr}\)。选一个满足 \(|\lambda_j|=L\) 的指数，再取 \(\theta=-\arg\lambda_j\)，则该方向的 \(H=L\)。前面的下界遂给
\[
 \limsup_{r\to\infty}\frac{\log M(r,F)}r=L.
\]
所以 \(F\) 的阶为一，型为 \(L\)。
\end{solution}

% 自拟：反向设计任意正阶和扩展型，允许任意稀疏非零系数支集并保留有限 Taylor 数据。
\Needspace{6\baselineskip}
\begin{exercise}[稀疏系数的阶型构造]\label{b3:ex:19-draft-02}
给定 \(\rho>0\)、\(\tau\in[0,\infty]\)、整数 \(N\ge0\)、次数不超过 \(N\) 的多项式 \(P\)，以及任意严格递增整数列
\[
 N<n_1<n_2<\cdots.
\]
构造超越整函数 \(f\)，使其阶为 \(\rho\)、型为 \(\tau\)，并且 \(f-P\) 的非零 Taylor 系数指标恰为 \(n_1,n_2,\ldots\)。构造应同时覆盖最小型、有限正型和最大型，并说明为什么对指标列不需要额外的密度条件。
\end{exercise}
\begin{solution}
按所需型取正数
\[
 s_k=
 \begin{cases}
 (\log(n_k+e))^{-1},&\tau=0,\\
 \tau,&0<\tau<\infty,\\
 \log(n_k+e),&\tau=\infty,
 \end{cases}
 \quad
 c_{n_k}=\left(\frac{e\rho s_k}{n_k}\right)^{n_k/\rho},
\]
并定义
\[
 f(z)=P(z)+\sum_{k=1}^\infty c_{n_k}z^{n_k}.
\]
其余高次系数都取零。三种情况下均有 \(s_k/n_k\to0\)，所以
\[
 |c_{n_k}|^{1/n_k}=\left(\frac{e\rho s_k}{n_k}\right)^{1/\rho}\longrightarrow0.
\]
全体系数的根也趋于零，级数在每个闭圆盘绝对一致收敛，定义整函数。所列高次系数全部为正且无穷多个非零，所以函数不是多项式。

又有 \(\log s_k=o(\log n_k)\)，从而
\[
 \frac{n_k\log n_k}{\log(1/|c_{n_k}|)}
 =\frac{\rho\log n_k}{\log n_k-\log(e\rho s_k)}
 \longrightarrow\rho.
\]
系数判据给出所需阶；有限多项式头部与零系数不改变这个上极限。型的判据则给
\[
 n_k|c_{n_k}|^{\rho/n_k}=e\rho s_k,\quad
 \tau(f)=\frac1{e\rho}\limsup_k e\rho s_k=\tau.
\]
当目标型为零或无穷时，最后的等式分别按趋零或趋于无穷解释。

整个计算只用到了 \(n_k\to\infty\)，没有用相邻指标的间距。因此无论这些指标多么稀疏，均可在每个非零系数上安排适当大小，保持指定的阶与型；任意给定的有限 Taylor 数据也保持不变。
\end{solution}

% 自拟：以双指数衰减系数求主导项的确切转折，再用最大项比较导出不同于正文高斯系数例子的零阶增长。
\Needspace{6\baselineskip}
\begin{exercise}[最大项的转折与零阶增长]\label{b3:ex:19-draft-03}
令
\[
 F(z)=\sum_{n=0}^\infty e^{-2^n}z^n,\quad
 \mu(r)=\max_{n\ge0}e^{-2^n}r^n.
\]
证明 \(F\) 为整函数，并精确求出最大项指标发生变化的半径。当 \(t=\log r>1\) 时，确定哪些整数 \(n\) 能使 \(nt-2^n=\log\mu(r)\)，包括并列最大情形。最后证明
\[
 \log M(r,F)=\frac{\log r\,\log\log r}{\log2}+O(\log r)
 \quad(r\to\infty),
\]
并据此说明 \(F\) 是零阶超越整函数，而且其最大模超过每个固定幂 \(r^d\)。
\end{exercise}
\begin{solution}
系数根为 \(e^{-2^n/n}\to0\)，所以收敛半径无穷。由于系数非负，三角不等式在正实轴取等号，故 \(M(r,F)=F(r)\)。

令 \(\phi_n(t)=nt-2^n\)。相邻项之差为
\[
 \phi_{n+1}(t)-\phi_n(t)=t-2^n,
\]
它随 \(n\) 严格递减。因此各项先增加后减少，\(n\ge1\) 为最大指标的充要条件是
\[
 2^{n-1}\le t\le2^n.
\]
若两个不等式均严格，最大指标唯一；若 \(t=2^n\)，恰有 \(n,n+1\) 两个并列指标。\(t<1\) 时指标为零，\(t=1\) 时指标零与一并列。所有转折半径因而正是
\(r=\exp(2^n)\)、\(n\ge0\)。

当 \(t>1\) 时，选择任一最大指标 \(n\)，有
\[
 n=\frac{\log t}{\log2}+O(1),\quad t\le2^n\le2t.
\]
所以
\[
 \log\mu(e^t)=nt-2^n=\frac{t\log t}{\log2}+O(t).
\]
最大项比较给
\(\mu(r)\le M(r,F)\le2\mu(2r)\)。将后一项的对数半径写成 \(t+\log2\)，上面的渐近式仍具有主项 \(t\log t/\log2\)，误差为 \(O(t)\)。上下夹逼便得所求最大模渐近。

于是
\[
 \frac{\log\log M(r,F)}{\log r}\longrightarrow0,
 \quad
 \frac{\log M(r,F)}{\log r}
 =\frac{\log\log r}{\log2}+O(1)\longrightarrow\infty.
\]
第一式给零阶，第二式说明对每个固定 \(d\)，最终 \(M(r,F)>r^d\)，并排除多项式。主导项虽然移动得很慢，仍会依次经过任意高的次数。
\end{solution}

% 自拟：将离散有限差分还原为插值多项式，再以整数零点和最小指数增长排除非零余项。
\Needspace{6\baselineskip}
\begin{exercise}[整数数据与多项式唯一性]\label{b3:ex:19-draft-04}
设整函数 \(f\) 满足：对每个 \(\varepsilon>0\)，存在 \(C_\varepsilon<\infty\)，使
\[
 |f(z)|\le C_\varepsilon e^{\varepsilon|z|}\quad(z\in\mathbb C).
\]
又设某个整数 \(m\ge1\) 使
\[
 \sum_{j=0}^m(-1)^{m-j}\binom mj f(n+j)=0
 \quad(n=0,1,2,\ldots).
\]
证明 \(f\) 是次数小于 \(m\) 的多项式，并用 \(f(0),\ldots,f(m-1)\) 写出它。构造一个阶一、有限正型的非多项式，满足相同的离散条件，从而说明增长假设的作用。
\end{exercise}
\begin{solution}
对整数数据记 \(\Delta f(n)=f(n+1)-f(n)\)，题设即 \(\Delta^mf(n)=0\)。定义多项式
\[
 P(z)=\sum_{k=0}^{m-1}(\Delta^kf)(0)\binom zk,
 \quad
 \binom z0=1,\quad
 \binom zk=\frac{z(z-1)\cdots(z-k+1)}{k!}\quad(k\ge1).
\]
其中每个 \((\Delta^kf)(0)\) 只使用 \(f(0),\ldots,f(k)\)。移位算子满足 \(E=I+\Delta\)，故对 \(0\le j<m\)，有限二项式展开给
\[
 f(j)=(E^jf)(0)=\sum_{k=0}^j\binom jk(\Delta^kf)(0)=P(j).
\]
多项式 \(P\) 的次数小于 \(m\)，因此 \(\Delta^mP=0\)。两列数据 \(f(n)\)、\(P(n)\) 满足同一个 \(m\) 阶递推式，最后一项的系数为一，且前 \(m\) 项相同。逐项递推便得 \(f(n)=P(n)\) 对所有非负整数成立。

令 \(h=f-P\)。多项式也满足任意正指数的增长界，因此若 \(h\not\equiv0\)，则
\(\log M(r,h)=o(r)\)。设原点重数为 \(\ell\)，写 \(h(z)=z^\ell g(z)\)、\(g(0)\ne0\)。所有正整数仍为 \(g\) 的零点。Jensen 的两半径估计于是给
\[
 \lfloor R\rfloor\log2
 \le\log M(2R,g)-\log|g(0)|
 =\log M(2R,h)-\ell\log(2R)-\log|g(0)|.
\]
右端是 \(o(R)\)，左端与 \(R\) 成正比，令 \(R\to\infty\) 即矛盾。所以 \(h=0\)，所求函数就是上面显式给出的 \(P\)。

函数 \(\sin\pi z\) 在所有整数上为零，故对任意 \(m\ge1\) 都满足题中的离散条件，却不是多项式。它的阶一、型为 \(\pi\)，不满足对每个 \(\varepsilon>0\) 的原增长界。可见把最小指数增长放宽为任意有限正常型，会失去该唯一性结论。
\end{solution}

% 对原自拟题实质改编：对同一组显式零点计算计数与增长，并在给定指数修正族内求精确最小型值，不使用分解定理。
\Needspace{6\baselineskip}
\begin{exercise}[同一组零点与可变的型]\label{b3:ex:19-draft-05}
对整数 \(m\ge1\) 和 \(a\in\mathbb C\)，令
\[
 F_a(z)=e^{az^m}(e^{z^m}-1).
\]
求其全部零点与重数、闭圆盘零点计数及收敛指数，并判定临界幂和是否收敛。证明 \(F_a\) 的阶为 \(m\)，型为
\[
 \tau(F_a)=\max\{|a|,|a+1|\}.
\]
确定这族函数的最小型值及达到该值的参数，并说明其圆周对数平均为什么不依赖 \(a\)。这里型的最小值只在给定的 \(F_a\) 族内比较。
\end{exercise}
\begin{solution}
指数前因子没有零点。方程 \(e^{z^m}=1\) 等价于 \(z^m=2\pi\mathrm i k\)、\(k\in\mathbb Z\)。\(k=0\) 给原点的 \(m\) 重零点，因为 \(e^{z^m}-1=z^m+O(z^{2m})\)。对 \(k\ne0\)，方程有 \(m\) 个互异根，且导数 \(mz^{m-1}e^{z^m}\) 在各根不为零，所以它们都是单零点。

于是对 \(r>0\)，闭圆盘计数为
\[
 n_{F_a}(r)=m+2m\left\lfloor\frac{r^m}{2\pi}\right\rfloor.
\]
全部非零零点的 \(s\) 次幂和为
\[
 2m\sum_{k=1}^\infty(2\pi k)^{-s/m}.
\]
它收敛当且仅当 \(s>m\)，故收敛指数是 \(m\)，临界 \(s=m\) 时发散。

记 \(H_a(w)=e^{(a+1)w}-e^{aw}\)，并令 \(L=\max\{|a|,|a+1|\}>0\)。上界为 \(|H_a(w)|\le2e^{L|w|}\)。在 \(a,a+1\) 中选一个模为 \(L\) 的数 \(\lambda\)，另一个记为 \(\nu\)。取 \(w=Re^{-\mathrm i\arg\lambda}\)，则
\[
 \operatorname{Re}(\lambda w)=LR,\quad
 \operatorname{Re}(\nu e^{-\mathrm i\arg\lambda})<L.
\]
后一不等式严格，因为 \(|\nu|\le L\)，而实部达到 \(L\) 会迫使 \(\nu=\lambda\)，与两指数相差一矛盾。因此存在 \(\delta>0\)，使
\[
 |H_a(w)|\ge e^{LR}(1-e^{-\delta R}).
\]
由 \(z\mapsto z^m\) 把半径 \(r\) 的圆周满射到半径 \(r^m\) 的圆周，得到
\[
 \log M(r,F_a)=Lr^m+O(1)
\]
于充分大的 \(r\)。所以阶为 \(m\)，型为 \(L\)。

三角不等式给
\[
 1=|(a+1)-a|\le|a+1|+|a|\le2L,
\]
故型的最小可能值为 \(1/2\)，在 \(a=-1/2\) 时达到，此时 \(F_a(z)=2\sinh(z^m/2)\)。若 \(L=1/2\)，上面两个不等式均取等号，\(a\) 与 \(a+1\) 必须位于同一直线的相反方向且模都为 \(1/2\)，故仍有 \(a=-1/2\)。此外取实数 \(a=L-1\)，可实现任意 \(L\ge1/2\)，这族函数的型恰遍历 \([1/2,\infty)\)。

最后，在每个圆周上有几乎处处的等式
\[
 \log|F_a(re^{\mathrm it})|
 =\operatorname{Re}(ar^me^{\mathrm imt})+\log|e^{r^me^{\mathrm imt}}-1|.
\]
第二项的零点奇性可积，第一项的平均为零，所以圆周对数平均完全不依赖 \(a\)。零点和这一平均保持不变，最大模的型却可以在上述范围内变化。
\end{solution}

% 对原自拟题实质改编：由导数界建立外域 Laurent 生成函数与积分重建，并用轮廓变形识别精确指数增长率。
\Needspace{6\baselineskip}
\begin{exercise}[导数生成函数与指数增长]\label{b3:ex:19-draft-06}
设 \(f\) 是非恒零整函数，且
\[
 T=\limsup_{r\to\infty}\frac{\log M(r,f)}r<\infty.
\]
证明
\[
 T=\limsup_{n\to\infty}|f^{(n)}(0)|^{1/n}.
\]
定义 Laurent 级数
\[
 \mathcal B_f(\zeta)=\sum_{n=0}^\infty\frac{f^{(n)}(0)}{\zeta^{n+1}}.
\]
证明它在 \(|\zeta|>T\) 全纯，且对任意 \(R>T\)、\(z\in\mathbb C\)，有
\[
 f(z)=\frac1{2\pi\mathrm i}\int_{|\zeta|=R}e^{z\zeta}\mathcal B_f(\zeta)\,\mathrm d\zeta.
\]
若 \(T>0\)，证明 \(\mathcal B_f\) 不可能全纯延拓到任何外域 \(|\zeta|>s\)、\(0\le s<T\)。
\end{exercise}
\begin{solution}
非恒零整函数的最大模在大圆上有一个正下界，因此 \(T\ge0\)。给定 \(t>T\)，有 \(M(r,f)\le C_t e^{tr}\)。对 \(n\ge1\)，Cauchy 估计并取 \(r=n/t\)，得到
\[
 |f^{(n)}(0)|\le n!C_t e^{tr}r^{-n}
 =C_t n!\left(\frac{et}{n}\right)^n.
\]
由 \(\log n!=n\log n-n+O(\log n)\)，有 \((n!)^{1/n}/n\to e^{-1}\)。所以若记右侧所求上极限为 \(L\)，便得 \(L\le t\)。令 \(t\downarrow T\)，得到 \(L\le T\)，也包括 \(T=0\) 的情形。

反向，给定 \(b>L\)，所有导数可用某个常数 \(A_b\) 统一估计为 \(|f^{(n)}(0)|\le A_b b^n\)，有限个初始指标由扩大常数吸收。Taylor 级数于是给
\[
 |f(z)|\le A_b\sum_{n=0}^\infty\frac{(b|z|)^n}{n!}=A_b e^{b|z|}.
\]
所以 \(T\le b\)，再令 \(b\downarrow L\) 即得 \(T=L\)。

这一系数根上极限说明 \(\mathcal B_f\) 在 \(|\zeta|>T\) 的每个紧集上绝对一致收敛，因而全纯。固定 \(R>T\)。Laurent 级数和 \(e^{z\zeta}\) 的幂级数在积分圆周上都绝对一致收敛，可以逐项相乘积分。只有指数为 \(\zeta^{-1}\) 的项留下，故
\[
 \frac1{2\pi\mathrm i}\int_{|\zeta|=R}e^{z\zeta}\mathcal B_f(\zeta)\,\mathrm d\zeta
 =\sum_{n=0}^\infty f^{(n)}(0)\frac{z^n}{n!}=f(z).
\]
这里 \(z\) 可以是任意固定复数，不需要限制在某个小圆盘。

若 \(T>0\) 且该函数全纯延拓到 \(|\zeta|>s\)、\(s<T\)，选择 \(s<R_1<T<R\)。被积函数在两圆之间的闭环域邻域全纯，Cauchy 定理允许将重建积分缩到 \(|\zeta|=R_1\)。于是对所有 \(z\)，有
\[
 |f(z)|\le R_1\max_{|\zeta|=R_1}|\mathcal B_f(\zeta)|\,e^{R_1|z|}.
\]
这给 \(T\le R_1<T\)，矛盾。因此精确指数增长率也限制了这个外域生成函数能够向内延拓的范围。整个论证容许 \(T=0\)；此时生成函数在任意非零 \(\zeta\) 处有定义，最后的严格缩小半径命题则没有需要讨论的非负 \(s\)。
\end{solution}

% 自拟：将多射线的多项式边界界转为全局幂界；分扇形选外点分母，保留临界最小型及非整数次数限制。
\Needspace{6\baselineskip}
\begin{exercise}[射线上的多项式控制]\label{b3:ex:19-draft-07}
取有限条互异射线，辐角依次为 \(\theta_1<\cdots<\theta_k<\theta_1+2\pi\)，其中 \(k\ge2\)，并令 \(\theta_{k+1}=\theta_1+2\pi\)、
\[
 \beta=\max_j(\theta_{j+1}-\theta_j),\quad \kappa=\frac\pi\beta.
\]
设整函数 \(f\) 在这些射线上满足
\[
 |f(re^{\mathrm i\theta_j})|\le C(1+r)^N
 \quad(r\ge0,\ 1\le j\le k),
\]
其中 \(N\ge0\) 为实数。若 \(\rho(f)<\kappa\)，或 \(\rho(f)=\kappa\) 且在该阶为最小型，证明 \(f\) 是次数不超过 \(\lfloor N\rfloor\) 的多项式。说明临界有限正型不能直接纳入此结论。
\end{exercise}
\begin{solution}
令 \(d=\lceil N\rceil\)。对相邻射线夹成的每个闭扇形 \(\overline S_j\)，取一个单位模的外点 \(a_j\notin\overline S_j\)。每个角宽都小于 \(2\pi\)，所以这样的外点存在。还有常数 \(c_j>0\)，使
\[
 |z-a_j|\ge c_j(1+|z|)\quad(z\in\overline S_j).
\]
事实上，在 \(|z|\ge2\) 时有 \(|z-a_j|\ge|z|-1\ge(1+|z|)/3\)；在剩下的紧集上，距离比为连续正函数，取其正最小值即可。

在 \(S_j\) 上定义
\[
 g_j(z)=\frac{f(z)}{(z-a_j)^d}.
\]
极点位于闭扇形之外，所以 \(g_j\) 在扇形全纯并连续到有限边界和顶点。两条边界射线上有
\[
 |g_j(z)|\le Cc_j^{-d}(1+|z|)^{N-d}\le Cc_j^{-d}.
\]
分母的下界还说明它保留了所需的增长上界。

若 \(\rho(f)<\kappa\)，取 \(\rho(f)<q<\kappa\)。每个扇形的临界指数 \(\pi/(\theta_{j+1}-\theta_j)\) 都不小于 \(\kappa\)，故可对 \(g_j\) 使用严格次临界 Phragmén--Lindelöf 原理。若 \(f\) 为阶 \(\kappa\) 的最小型，则在角宽等于 \(\beta\) 的扇形中使用临界最小型版本；在较窄扇形中，\(\kappa\) 已严格小于该扇形的临界指数，仍可使用严格版本。两种情况下均得 \(g_j\) 在其扇形有界。

因此存在 \(A_j<\infty\)，使 \(|f(z)|\le A_j(1+|z|)^d\) 于 \(\overline S_j\)。扇形只有有限个，取这些常数的最大值便得到全平面的同类幂界。多项式增长判别说明 \(f\) 是次数不超过 \(d\) 的多项式。若其实际次数 \(\ell>N\)，沿任一固定射线都有
\[
 |f(re^{\mathrm i\theta_j})|=|c_\ell|r^\ell(1+o(1)),
\]
与原来的 \((1+r)^N\) 界矛盾，所以实际次数至多 \(\lfloor N\rfloor\)。零多项式自然包括在结论中。这里每个商只用于自己的扇形，不要求把带外部极点的各个商合成一个整函数。

为看清临界限制，取整数 \(m\ge1\) 及射线 \(\theta_j=j\pi/m\)、\(j=0,\ldots,2m-1\)。它们的最大角宽为 \(\pi/m\)，临界指数为 \(m\)。整函数 \(e^{\mathrm iz^m}\) 在每条射线上模恒为一，阶 \(m\)、型一，却不是多项式。因此临界处仅有有限正型不足以替代最小型。
\end{solution}

% 自拟：为两侧高斯模界构造精确的三次调和修正，给出尖锐内部估计及机械同高插值的反例。
\Needspace{6\baselineskip}
\begin{exercise}[高斯边界权的三线比较]\label{b3:ex:19-draft-08}
设 \(a_0,a_1\ge0\)，\(f\) 在条带 \(0<\operatorname{Re}z<1\) 全纯，并连续到闭条带。假设存在 \(C>0\)、\(0<q<\pi\)，使
\[
 \log^+|f(x+\mathrm iy)|\le Ce^{q|y|}\quad(0\le x\le1),
\]
并且
\[
 |f(\mathrm iy)|\le e^{a_0y^2},\quad
 |f(1+\mathrm iy)|\le e^{a_1y^2}\quad(y\in\mathbb R).
\]
构造次数不超过三的全纯多项式 \(Q\)，使其在两条边界上的实部分别为 \(a_0y^2\)、\(a_1y^2\)，由此求出逐点最优的内部模界，并确定在一个内部点取等号的全部函数。最后用 \(a_0=a_1>0\) 说明，直接把同一高度的两个边界界作几何插值一般不正确。
\end{exercise}
\begin{solution}
令 \(\delta=a_1-a_0\)，取
\[
 Q(z)=-\frac\delta3z^3-a_0z^2+\frac{2a_0+a_1}{3}z.
\]
由 \(\operatorname{Re}z^2=x^2-y^2\) 和 \(\operatorname{Re}z^3=x^3-3xy^2\)，可算得
\[
 \begin{aligned}
 H(x,y)\coloneqq\operatorname{Re}Q(x+\mathrm iy)
 &=\bigl((1-x)a_0+xa_1\bigr)y^2\\
 &\quad+\frac{x(1-x)}3\bigl((2-x)a_0+(1+x)a_1\bigr).
 \end{aligned}
\]
这给 \(H(0,y)=a_0y^2\)、\(H(1,y)=a_1y^2\)。额外的只依赖 \(x\) 的项使整个表达式调和；只对两侧 \(y^2\) 系数作线性插值，通常没有这个性质。

令 \(g=f e^{-Q}\)。它在两竖边的模不超过一。由于 \(x\in[0,1]\)，上式还给 \(|\operatorname{Re}Q(x+\mathrm iy)|\le C_1(1+y^2)\)，所以
\[
 \log^+|g(x+\mathrm iy)|
 \le Ce^{q|y|}+C_1(1+y^2)
 \le C_2e^{q|y|}.
\]
最后一步使用 \(q>0\)，常数可以依赖于 \(q,a_0,a_1\)。因此适用本节允许竖向增长的三线比较，两条边界常数都为一，得到 \(|g|\le1\)。所求内部界为
\[
 \boxed{\ |f(x+\mathrm iy)|\le
 \exp\left(\bigl((1-x)a_0+xa_1\bigr)y^2
 +\frac{x(1-x)}3\bigl((2-x)a_0+(1+x)a_1\bigr)\right).\ }
\]

对任意实数 \(\alpha\)，函数 \(f(z)=e^{\mathrm i\alpha}e^{Q(z)}\) 在两侧和内部每一点都取等号。它的 \(\log^+\) 在闭条带上为 \(O(1+y^2)\)，满足题设增长条件，所以这个界逐点最优。反之，若任意可行 \(f\) 在一个内部点取等号，则 \(g\) 在内部取得最大模一，最大模原理给 \(g\equiv e^{\mathrm i\alpha}\)。因此上面的函数也就是全部等号情形。

若 \(a_0=a_1=a>0\)，则 \(Q(z)=az(1-z)\)，比较函数的实部简化为
\[
 H(x,y)=ay^2+ax(1-x).
\]
特别地，\(f(z)=e^{az(1-z)}\) 在两竖边均有模 \(e^{ay^2}\)，在 \(z=1/2\) 却有模 \(e^{a/4}>1\)。同一高度 \(y=0\) 的两个边界界都是一，其几何插值仍为一，遗漏了内部必要的 \(ax(1-x)\) 项。三线定理中的边界常数必须对整条边界成立，不能逐个高度代入不同的常数后仍沿用原公式。
\end{solution}
